\begin{figure}[H]
    \centering    \includegraphics[width=0.8\textwidth]{Figures}
    \caption{This figure...}
    \label{fig:}
\end{figure}

$
e^{\mathcal{I}k \Delta x}e^{\mathcal{I}l \Delta y}+e^{\mathcal{I}k \Delta x}-e^{\mathcal{I}l \Delta y}-1  \\
+e^{\mathcal{I}k \Delta x}+ e^{\mathcal{I}k \Delta x}e^{-\mathcal{I}l \Delta y}-1-e^{\mathcal{I}l \Delta y} \\
-e^{\mathcal{I} \Delta y}-1+e^{-\mathcal{I}k \Delta x}e^{\mathcal{I}l \Delta y}+e^{-\mathcal{I}k \Delta x} \\ 
-1-e^{-\mathcal{I}l \Delta y}+e^{-\mathcal{I}k \Delta x}+e^{-\mathcal{I}k \Delta x}e^{-\mathcal{I}l \Delta y} \\
$

$
e^{\mathcal{I}k \Delta x}e^{\mathcal{I}l \Delta y}+e^{\mathcal{I}l \Delta y}-e^{\mathcal{I}k \Delta x}-1  \\
+e^{\mathcal{I}l \Delta y}+ e^{-\mathcal{I}k \Delta x}e^{\mathcal{I}l \Delta y}-1-e^{\mathcal{I}k \Delta x} \\
-e^{\mathcal{I}k \Delta x}-1+e^{\mathcal{I}k \Delta x}e^{-\mathcal{I}l \Delta y}+e^{-\mathcal{I}l \Delta y} \\ 
-1-e^{-\mathcal{I}k \Delta x}+e^{-\mathcal{I}l \Delta y}+e^{-\mathcal{I}k \Delta x}e^{-\mathcal{I}l \Delta y} \\
$

If we now assume that $\Delta x = \Delta y$, which implies that $\mu_x = \mu_y = \mu$, and that $k=l$ we can simplify the above expression to get
\begin{align*}
    \left(\lambda - \lambda^{-1}\right)^2 &= 2\mu^2\left[ 
    -e^{\mathcal{I}k \Delta x}+e^{-\mathcal{I}k \Delta x} +e^{2\mathcal{I}k \Delta x} 
     +e^{-2\mathcal{I}k \Delta x} -2 \right] \\
     &= 4\mu^2\left[ -\mathcal{I}\sin(k\Delta x)+\cos(2k\Delta x) - 1 \right] \\
     &= -4\mu^2\left[\sin^2(k\Delta x) + \mathcal{I}\sin(k\Delta x) \right]
\end{align*}



#######################

\section{The stability of the non-rotating 2D SWE  on a B-grid}

\textbf{(a)} The non-rotating 2-D SWE equations are given as 
\begin{subequations}
\begin{align}
\frac{\partial u}{\partial t}& = -g \frac{\partial h}{\partial x} \\
\frac{\partial v}{\partial t}& = -g \frac{\partial h}{\partial y} \\
\frac{\partial h}{\partial t}& + H \left( \frac{\partial u}{\partial x} + \frac{\partial v}{\partial y} \right) = 0.
\end{align}
\label{eq:SWE}
\end{subequations}
Discretising the equations on an Arakawa B-grid by applying the leapfrog scheme one obtains
\begin{subequations}
\begin{align}
\frac{u_{i,j}^{n+1} - u_{i,j}^{n-1}}{2\Delta t} &= -g \frac{h_{i+1,j+1}^{n} + h_{i+1,j}^{n}-h_{i,j+1}^{n}-h_{i,j}^{n}}{2\Delta x} \\
\frac{v_{i,j}^{n+1} - v_{i,j}^{n-1}}{2\Delta t} &= -g \frac{h_{i,j+1}^{n} + h_{i+1,j+1}^{n}-h_{i+1,j}^{n}-h_{i,j}^{n}}{2\Delta y}
\\
\frac{h_{i,j}^{n+1} - h_{i,j}^{n-1}}{2\Delta t}& = - H \left( \frac{u_{i,j}^{n} + u_{i,j-1}^{n}-u_{i-1,j}^{n}-u_{i-1,j-1}^{n}}{2\Delta x} + \frac{v_{i,j}^{n} + v_{i-1,j}^{n}-v_{i,j-1}^{n}-v_{i-1,j-1}^{n}}{2\Delta y} \right).
\end{align}
\label{eq:discretised_SWE}
\end{subequations}

\textbf{(b)} To solve for the stability criterion we apply a von Neumann stability analysis. We assume that the solution can be written as a Fourier mode of the form
\begin{equation}\begin{pmatrix}
u_{i,j}^{n} \\
v_{i,j}^{n} \\
h_{i,j}^{n}
\end{pmatrix} = \begin{pmatrix}u_0 \\
v_0 \\
h_0
\end{pmatrix} e^{\mathcal{I}(k i \Delta x + l j \Delta y - \omega n \Delta t)},
\end{equation}
where $k$ and $l$ are the wavenumbers in the $x$ and $y$ direction, respectively, $\omega$ is the frequency and $\mathcal{I}$ is the imaginary unit. Substituting this into the spatial and temporal time steps one can transform \autoref{eq:discretised_SWE} to 
\begin{subequations}
\begin{align}
\frac{\lambda u_0 - \lambda^{-1} u_0}{2\Delta t} &= -g \frac{h_0e^{\mathcal{I}k \Delta x}e^{\mathcal{I}l \Delta y}+ h_0e^{\mathcal{I}k\Delta x} - h_0e^{\mathcal{I}l \Delta y}-h_0}{2\Delta x} \\
\frac{\lambda v_0 - \lambda^{-1} v_0}{2\Delta t} &= -g \frac{h_0e^{\mathcal{I}k \Delta x}e^{\mathcal{I}l \Delta y}+ h_0e^{\mathcal{I}l\Delta y} - h_0e^{\mathcal{I}k \Delta x}-h_0}{2\Delta y} \\
\begin{split}
\frac{\lambda h_0 - \lambda^{-1} h_0}{2\Delta t} = - H  &\left( \frac{u_0+u_0e^{-\mathcal{I}l \Delta y} - u_0e^{-\mathcal{I}k \Delta x}- u_0e^{-\mathcal{I}k \Delta x}e^{-\mathcal{I}l \Delta y}}{2\Delta x} \right. \\
&\quad \left. + \frac{v_0+v_0e^{-\mathcal{I}k \Delta x} - v_0e^{-\mathcal{I}l \Delta y}- v_0e^{-\mathcal{I}k \Delta x}e^{-\mathcal{I}l \Delta y}}{2\Delta y} \right).
\end{split}
\end{align}
\label{eq:discretised_SWE_Fourier}
\end{subequations}
Simplifying the above equations above gives
\begin{subequations}
\begin{align*}
u_0 &= -\frac{1}{\lambda - \lambda^{-1}} \frac{g h_0 \Delta t}{\Delta x} \left(e^{\mathcal{I}k \Delta x}e^{\mathcal{I}l \Delta y}+ e^{\mathcal{I}k\Delta x} - e^{\mathcal{I}l \Delta y}-1\right) \\
v_0 &= -\frac{1}{\lambda - \lambda^{-1}} \frac{gh_0\Delta t}{\Delta y} \left(e^{\mathcal{I}k \Delta x}e^{\mathcal{I}l \Delta y}+ e^{\mathcal{I}l\Delta y} - e^{\mathcal{I}k \Delta x}-1\right) \\
\begin{split}
    \left(\lambda - \lambda^{-1}\right) h_0 = &- \frac{Hu_0\Delta t}{\Delta x}  \left(1+e^{-\mathcal{I}l \Delta y} - e^{-\mathcal{I}k \Delta x}- e^{-\mathcal{I}k \Delta x}e^{-\mathcal{I}l \Delta y}\right)  \\ 
    &- \frac{Hv_0\Delta t}{\Delta y}\left(1+e^{-\mathcal{I}k \Delta x} - e^{-\mathcal{I}l \Delta y}- e^{-\mathcal{I}k \Delta x}e^{-\mathcal{I}l \Delta y}\right).
\end{split}
\end{align*}
\end{subequations}
Now inserting the expressions for $u_0$ and $v_0$ into the last equation gives
\begin{equation*}
    \begin{split}
    \left(\lambda - \lambda^{-1}\right)^2 &= \frac{gH\left(\Delta t\right)^2}{\left(\Delta x\right)^2} \left(e^{\mathcal{I}k \Delta x}e^{\mathcal{I}l \Delta y}+ e^{\mathcal{I}k\Delta x} - e^{\mathcal{I}l \Delta y}-1\right)\left(1+e^{-\mathcal{I}l \Delta y} - e^{-\mathcal{I}k \Delta x}- e^{-\mathcal{I}k \Delta x}e^{-\mathcal{I}l \Delta y}\right) \\ 
    &
    + \frac{gH\left(\Delta t\right)^2}{\left(\Delta y\right)^2} \left(e^{\mathcal{I}k \Delta x}e^{\mathcal{I}l \Delta y}+ e^{\mathcal{I}l\Delta y} - e^{\mathcal{I}k \Delta x}-1\right)\left(1+e^{-\mathcal{I}k \Delta x} - e^{-\mathcal{I}l \Delta y}- e^{-\mathcal{I}k \Delta x}e^{-\mathcal{I}l \Delta y}\right).
    \end{split}
\end{equation*}
If we now use the Courant number $\mu_x = \tfrac{\sqrt{gH}\Delta t}{\Delta x}$ and $\mu_y = \tfrac{\sqrt{gH}\Delta t}{\Delta y}$ we can expand the above expression to get
\begin{equation*}
    \begin{split}
    \left(\lambda - \lambda^{-1}\right)^2 = &\mu_x^2\left[  e^{\mathcal{I}k \Delta x}e^{\mathcal{I}l \Delta y}+e^{\mathcal{I}k \Delta x}-e^{\mathcal{I}l \Delta y}-1 \right. \\
    &\quad+e^{\mathcal{I}k \Delta x}+ e^{\mathcal{I}k \Delta x}e^{-\mathcal{I}l \Delta y}-1-e^{\mathcal{I}l \Delta y} \\
     &\quad-e^{\mathcal{I}l \Delta y}-1+e^{-\mathcal{I}k \Delta x}e^{\mathcal{I}l \Delta y}+e^{-\mathcal{I}k \Delta x} \\ 
     &\quad-1-e^{-\mathcal{I}l \Delta y}+e^{-\mathcal{I}k \Delta x}+e^{-\mathcal{I}k \Delta x}e^{-\mathcal{I}l \Delta y} 
    \left.\right] \\
    \vspace{1cm}
    +&\mu_y^2\left[ e^{\mathcal{I}k \Delta x}e^{\mathcal{I}l \Delta y}+e^{\mathcal{I}l \Delta y}-e^{\mathcal{I}k \Delta x}-1  \right.\\
    &\quad+e^{\mathcal{I}l \Delta y}+ e^{-\mathcal{I}k \Delta x}e^{\mathcal{I}l \Delta y}-1-e^{\mathcal{I}k \Delta x} \\
    &\quad-e^{\mathcal{I}k \Delta x}-1+e^{\mathcal{I}k \Delta x}e^{-\mathcal{I}l \Delta y}+e^{-\mathcal{I}l \Delta y} \\ 
    &\quad\left.-1-e^{-\mathcal{I}k \Delta x}+e^{-\mathcal{I}l \Delta y}+e^{-\mathcal{I}k \Delta x}e^{-\mathcal{I}l \Delta y} \right]. 
    \end{split}
\end{equation*}
If we now assume that $\Delta x = \Delta y$, which implies that $\mu_x = \mu_y = \mu$, we thus simplify the above expression to get
\begin{equation*}
    \begin{split}
    \left(\lambda - \lambda^{-1}\right)^2 = &\mu^2\left[  e^{\mathcal{I}k \Delta x}e^{\mathcal{I}l \Delta y}+e^{\mathcal{I}k \Delta x}-e^{\mathcal{I}l \Delta y}-1 \right. \\
    &\quad+e^{\mathcal{I}k \Delta x}+ e^{\mathcal{I}k \Delta x}e^{-\mathcal{I}l \Delta y}-1-e^{\mathcal{I}l \Delta y} \\
     &\quad-e^{\mathcal{I}l \Delta y}-1+e^{-\mathcal{I}k \Delta x}e^{\mathcal{I}l \Delta y}+e^{-\mathcal{I}k \Delta x} \\ 
     &\quad-1-e^{-\mathcal{I}l \Delta y}+e^{-\mathcal{I}k \Delta x}+e^{-\mathcal{I}k \Delta x}e^{-\mathcal{I}l \Delta y} 
    \left.\right] \\
    \vspace{1cm}
    +&\mu^2\left[\right]
    \end{split}
\end{equation*}

